\section{Metodologia e arquitetura}
O pacote separa modelo de planta, métodos de análise, validação e
aplicação. O núcleo usa Python 3.12 e NumPy. O sistema define
coeficientes em ordem decrescente e avalia a razão por Horner, com
uso de polinômios reversos para $|s|>1$. A instrução de teste verifica
as importações permitidas no núcleo. SciPy, Python-control, SymPy e
Mpmath são executados na camada de validação.

\subsection{Contorno numérico e margens}
A malha contém frequências positivas e negativas, os arcos de desvio
e o arco externo. Para $A$ mônico, o código usa
$1-\sum_{i\ge1}|a_i|R^{-i}$ como cota inferior normalizada do
denominador no exterior, quando positiva. A cota superior da razão
é calculada sobre a mesma escala. Para um modelo estritamente
próprio e com atraso, $|e^{-s\tau}|\le1$ no semiplano direito permite
conter a cauda em um disco que exclui $-1$. Sem atraso, a cota se
aplica à diferença em relação ao ganho de passagem direta. Com atraso,
um ganho de passagem direta de módulo maior ou igual a um fica fora
da certificação desse contorno.

A contagem exige duas malhas sucessivas com o mesmo inteiro,
incremento angular máximo inferior a $\pi/3$ e distância do polígono
ao ponto crítico acima de $10^{-7}$. A distância é calculada nos
segmentos e nos vértices. O refinamento local bissecta segmentos cujo
mapeamento varia rapidamente ou apresenta curvatura relevante em relação
à distância ao ponto crítico. Cruzamentos de ganho unitário são pontos
adicionais da malha. A resposta próxima a $-1$ em um cruzamento
recebe o estado de fronteira. Um resultado sem resolução suficiente
recebe diagnóstico indeterminado. Quando a distância de um polo aberto ao eixo é numericamente
ambígua, a contagem fica indeterminada. Os raios de indentação
são reduzidos para preservar polos fechados próximos ao eixo.
Em sistemas racionais, a distância às raízes de $A+KB$ fornece
essa escala. Em sistemas atrasados, verifica-se a dominância do
termo de realimentação sobre o denominador em torno do polo.
A cota externa limita a região de interesse. A resolução no
interior continua empírica.

Para o ganho unitário, calcula-se o polinômio em $x=\omega^2$
a partir de $|KB(\imagunit\omega)|^2-|A(\imagunit\omega)|^2$. Cada candidato
encontrado é verificado na resposta complexa. Para a fase racional,
o produto $KB(\imagunit\omega)\overline{A(\imagunit\omega)}$ identifica a
interseção com o eixo real. Com atraso, a busca por semieixo negativo
ocorre em uma banda finita explícita. As passagens de Bode são
registradas em rad/s e as margens em razão de ganho e graus.

\subsection{Métodos polinomiais e resposta ao degrau}
Routh calcula a recorrência em ponto flutuante. Pivôs pequenos e
linhas nulas acionam uma alternativa racional para reduzir cancelamento.
A contagem no eixo usa as raízes nulas e os fatores comuns das partes
par e ímpar do polinômio. O lugar das raízes calcula as raízes
do polinômio para cada $K$ e conecta as amostras por atribuição
de custo mínimo sobre distâncias no plano complexo. O cálculo de
cruzamentos verifica os candidatos do produto de polinômios em
$\imagunit\omega$. Em caso de queda da ordem, a trajetória pode possuir
um ramo no infinito. O degrau usa uma realização companheira e uma
exponencial matricial calculada por série e quadraturas sucessivas.
Para plantas atrasadas, o degrau, Routh e lugar das raízes usam a
mesma aproximação de Padé escolhida pelo usuário.

\subsection{Plantas e hipóteses}
A Tabela~\ref{tab:systems} resume as plantas. Em A, os polos são
$-1,-2,-3$ e a fronteira de ganho é $K=10$ em $\sqrt{11}$ rad/s.
Em B, os polos são $1$ e $-2$. O modelo representa uma dinâmica
linearizada com rigidez efetiva negativa, como um pêndulo
invertido amortecido em torno do equilíbrio vertical. A escrita
normalizada é $\ddot x+\dot x-2x=u$. A fronteira em $K=2$ é
um cruzamento por frequência zero. Em C, os polos são $-1,-2$
e o zero é $+1$. O degrau inicia com inclinação negativa.
O polinômio $s^2+(3-K)s+(2+K)$ prevê cruzamento em $K=3$ e
$\sqrt5$ rad/s. Em D, o polo racional é $-1$ e a
primeira fronteira de atraso no ganho nominal é
$\tau_*=2\pi/(3\sqrt3)$ s, obtida de \eqref{eq:delay}.
As plantas A, B e D não têm zeros finitos.

\begin{table}[!t]
\centering
\caption{Plantas de teste e parâmetros nominais.}
\label{tab:systems}
\begin{tabular}{clcc}
\toprule
Caso & $B(s)/A(s)$ & $K$ & $\tau$ (s)\\
\midrule
A & $6/(s^3+6s^2+11s+6)$ & 1 & 0\\
B & $1/(s^2+s-2)$ & 4 & 0\\
C & $(1-s)/(s^2+3s+2)$ & 1 & 0\\
D & $2/(s+1)$ & 1 & 0,8\\
\bottomrule
\end{tabular}
\end{table}

\subsection{Validação e medições}
Os comparadores usam resultados analíticos onde há uma fronteira
conhecida e métodos de bibliotecas fora do núcleo. Assim é possível
registrar diferenças entre a representação em coeficientes e a avaliação
da função, a contagem de raízes e as convenções sobre cruzamento em DC.
Para os benchmarks, CPU e tempo de parede são medidos após aquecimento.
A carga de módulos, a geração de gráficos e a escrita em disco ficam
fora da seção cronometrada. Dados brutos guardam cada repetição e a
descrição do ambiente. A memória é medida à parte, sobre alocações
Python. A sensibilidade em raízes simples é interpretada pela relação
de primeira ordem $|\delta r|\simeq |\delta C(r)|/|C'(r)|$,
com distinção entre erro frente à construção analítica e resíduo
escalado do polinômio armazenado \cite{wilkinson1959}.
